% =====================================================================
%  Partie B — Questions à réponse courte
%
%  Modèle :
%    \begin{question}[points=3, lines=6]  % lines : espace dans le cahier
%    \begin{question}[points=3, box=6cm]  % cadre vide (ex. écrire du code)
%      Énoncé : verbe d'action du niveau taxonomique visé (Bloom),
%      qui indique le type de réponse attendu (Nommez, Expliquez,
%      Déterminez, Comparez, Écrivez…)
%      \begin{solution}
%        Réponse attendue (corrigé seulement)
%      \end{solution}
%    \end{question}
%
%  Question en plusieurs parties a), b)… : chaque sous-question a son
%  propre espace dans le cahier, précédé de sa lettre (les options
%  lines= et box= de la question sont alors ignorées) :
%    \begin{question}[points=5]
%      Énoncé commun
%      \begin{subquestion}[lines=3]
%        Expliquez…
%        \begin{solution} ... \end{solution}
%      \end{subquestion}
%      \begin{subquestion}[box=5cm]
%        Écrivez…
%        \begin{solution} ... \end{solution}
%      \end{subquestion}
%    \end{question}
%
%  Code tiré d'un fichier :
%    \codefile{java}{questions/code/Compte.java}
%    \codefile[firstline=8, lastline=13]{java}{questions/code/Compte.java}
% =====================================================================

\exampart{Questions à réponse courte}

\begin{instructions}
  Répondez dans l'espace prévu du cahier de réponses.
\end{instructions}

% --- Java, depuis un fichier, en plusieurs parties --------------------
% Chaque sous-question a son espace dans le cahier : lignes ou cadre.
\begin{question}[points=5]
  L'exécution de la méthode \codeinline{main} de la classe suivante
  provoque une erreur.
  \codefile{java}{questions/code/Compte.java}
  \begin{subquestion}[lines=3]
    Expliquez cette erreur dans vos mots.
    \begin{solution}
      Le second retrait (80\,\$) dépasse le solde (70\,\$) :
      une \codeinline{IllegalArgumentException} « Solde insuffisant »
      est levée et la ligne 23 n'est jamais exécutée.
    \end{solution}
  \end{subquestion}
  \begin{subquestion}[lines=2]
    Indiquez la ligne où l'erreur est déclenchée.
    \begin{solution}
      L'exception est levée à la ligne 10, lors de l'appel de la ligne 22.
    \end{solution}
  \end{subquestion}
  \begin{subquestion}[box=5cm]
    Réécrivez les lignes 21 à 23 pour que le programme affiche
    « Retrait refusé » au lieu de s'arrêter, puis affiche le solde.
    \begin{solution}
\begin{code}[numbers=none]{java}
try {
    c.retirer(30);
    c.retirer(80);
} catch (IllegalArgumentException e) {
    System.out.println("Retrait refusé");
}
System.out.println(c.getSolde());
\end{code}
    \end{solution}
  \end{subquestion}
\end{question}

% --- HTML + CSS : écrire du code -------------------------------------
\begin{question}[points=2, box=6cm]
  Voici le menu d'une page Web.
\begin{code}{html}
<!DOCTYPE html>
<html lang="fr">
<head>
  <meta charset="utf-8">
  <title>Menu</title>
  <link rel="stylesheet" href="style.css">
</head>
<body>
  <nav class="menu">
    <a href="index.html">Accueil</a>
    <a href="cours.html">Cours</a>
    <a href="contact.html">Contact</a>
  </nav>
</body>
</html>
\end{code}
  Écrivez les règles CSS qui affichent les liens du menu sur une même
  ligne, espacés de \codeinline{1rem}, sans soulignement.
  \begin{solution}
\begin{code}[numbers=none]{css}
.menu {
  display: flex;
  gap: 1rem;
}
.menu a {
  text-decoration: none;
}
\end{code}
  \end{solution}
\end{question}

% --- Node.js, depuis un fichier --------------------------------------
\begin{question}[points=2, lines=4]
  Le serveur Node.js suivant est démarré.
  \codefile{nodejs}{questions/code/serveur.js}
  Déterminez le code d'état HTTP et le texte que reçoit le navigateur
  pour une requête vers \codeinline{/cours}.
  \begin{solution}
    Code 404, texte « Introuvable » (l'URL n'est pas \codeinline{/}).
  \end{solution}
\end{question}

% --- Python : écrire une fonction ------------------------------------
\begin{question}[points=4, box=7cm]
  Écrivez une fonction Python \codeinline{moyenne(notes)} qui retourne la
  moyenne d'une liste de nombres. Si la liste est vide, la fonction
  retourne \codeinline{None}.
  \begin{solution}
\begin{code}[numbers=none]{python}
def moyenne(notes):
    if len(notes) == 0:
        return None
    return sum(notes) / len(notes)
\end{code}
    Critères : cas de la liste vide (1 pt), calcul correct (2 pts),
    valeur retournée et non affichée (1 pt).
  \end{solution}
\end{question}
